% secV3_data.tex -- Section 3. CONTAINER (27 Oct 2026): text to be written by the author.
% Word budget: 550 words (JPART count: running text, footnotes and citations as printed).
% Outline, numbers and sources: referee_EE_2026-10/ESQUELETO.md, section 3. The archived text
% (_arquivo_texto_IA_2026-10-27/) is not a source.

\section{Data}\label{sec:data}

We combine child-level administrative records provided by the SME with aggregate administrative publications, school enrolment data, and demographic records. We distinguish the \emph{request record}, which tracks individual requests, from the \emph{published waiting list}, which reports administrative counts at specific dates. Neither measures all families potentially interested in childcare.

The application files for 2010--2018 contain 3,298,064 records, including child identifiers, application dates, location, and dated referral and enrolment events observed through November 2019. We retain the first valid application per child, yielding 1,524,560 children (Table~\ref{tab:T1_fila}). Our unit of analysis for application flows is thus the child's first \emph{request}, not each administrative record. Because coverage is incomplete in the earliest years, analyses of entries and subsequent trajectories focus on first requests made from July 2012 onward. A \emph{referral} is an offer or assignment of a place, distinct from an effective \emph{enrolment}; children may also exit the childcare list upon reaching 47 months. Application timestamps are recorded to the second, although 2.7\% of first requests share a recorded second with another request (Table~\ref{tab:N1}). A guardian's CPF is recorded for 98.1\% of these children; we treat the identifier as present when it is non-missing and positive, without verifying its check digit. Separately provided Escola On-Line (EOL) enrolment records, linked by child identifier, permit comparisons with enrolment information recorded outside the application extract. Neither child-level source is publicly available. A separate capacity panel covers 2,734 centres across 13 semesters, with 86,982 centre--semester--age-group observations on authorized places.

We reconstruct the published series from the SME's original open-data files for 2006--2016 and subsequent administrative spreadsheets \citep{sme_dados_abertos,sme_demanda_escolar}. Municipal Law 14,127/2006 and Decree 47,155/2006 mandate quarterly publication of demand statistics by district and educational stage \citep{saopaulo2006lei14127}. SME Technical Note 01/2017 \citep{sme2017nt01} defines \emph{demanda} as registered applications without a referral, excluding cases with a referral, withdrawal or refusal, expired requests, school-transfer applications, and residents of other municipalities. The published figures are snapshots of the administrative system, rather than measures of unexpressed demand. Our assembled series contains 89 snapshots, quarterly through 2022 and bimonthly from 2023. Table~\ref{tab:series} reports the quarterly counts through 2022 and the June and December counts from 2023. Among 43 corresponding PDF releases checked against open-data totals, 37 match; the remaining six have formats or disaggregations that preclude a direct comparison. We use the underlying administrative tabulations.

For population coverage, the May School Census provides childcare-stage enrolment across all school networks and a municipal-and-partnered subset \citep{inep_censo_escolar_c}. For registered-demand calculations, we compare the SME's December enrolment series with the municipal-and-partnered School Census count; their reference dates and coverage differ. Population inputs comprise district-level live births from SINASC, the 2010 and 2022 Demographic Censuses, and SEADE's published age-specific projections \citep{sms_sinasc_tabnet,ibge_sidra1378,ibge_sidra9514,seade_idade_escolar}. Our principal under-four population series scales birth cohorts to both Census benchmarks. Alternatives interpolate directly between the two Census counts or between SEADE projections published for 2005, 2010, 2015, 2020, and 2025. Section~\ref{sec:goal} defines these constructions and the resulting compliance measures.

The application records were also used in \citet{soares2021ecd}, which studies childcare expansion and child outcomes. Our research question is different: we do not identify the effect of a programme, but examine whether a statutory performance claim is reproducible.
